% !TeX program = xelatex
\documentclass[12pt,a4paper]{article}
\usepackage[a4paper,margin=2.3cm]{geometry}
\usepackage{xcolor}
\usepackage{booktabs,longtable,array}
\usepackage{enumitem}
\usepackage{listings}
\usepackage[unicode,hidelinks]{hyperref}
\usepackage{xepersian}

\IfFontExistsTF{Amiri}{\settextfont{Amiri}}{\settextfont{FreeSerif}}
\IfFontExistsTF{TeX Gyre Termes}{\setlatintextfont{TeX Gyre Termes}}{\setlatintextfont{Latin Modern Roman}}

\setlength{\parindent}{1cm}
\setlength{\parskip}{2pt}
\linespread{1.25}
\setlist{itemsep=2pt,topsep=4pt}
\setlist[itemize]{label={\ensuremath{\bullet}}}
\definecolor{guideblue}{RGB}{16,64,120}
\hypersetup{pdftitle={راهنمای قالب پروژه کارشناسی دانشگاه بناب},pdfauthor={محمد امیر بابازاد}}

\lstset{
  basicstyle=\ttfamily\small,
  frame=single,
  backgroundcolor=\color{black!3},
  rulecolor=\color{black!20},
  breaklines=true,
  showstringspaces=false
}

\newcommand{\filepath}[1]{\lr{\nolinkurl{#1}}}
\newcommand{\commandname}[1]{\lr{\texttt{\textbackslash#1}}}

\title{\textbf{دفترچه راهنمای قالب پروژه کارشناسی دانشگاه بناب}}
\author{ویژه پروژه کارشناسی علوم کامپیوتر\\محمد امیر بابازاد}
\date{نسخه ۱٫۰ --- شهریور ۱۴۰۵}

\begin{document}
\maketitle
\thispagestyle{empty}
\vfill
\begin{center}
این راهنما روش نصب، اجرا، شخصی‌سازی و رفع خطاهای متداول قالب را توضیح می‌دهد.
\end{center}
\vfill
\newpage
\tableofcontents
\newpage

\section{معرفی}
این قالب برای گزارش پروژه کارشناسی رشته علوم کامپیوتر دانشگاه بناب آماده شده است. قواعد صفحه‌آرایی از دستورالعمل ارسالی دانشگاه استخراج شده‌اند و موارد مخصوص پایان‌نامه تحصیلات تکمیلی، مانند فرم دفاع کارشناسی ارشد، به‌طور پیش‌فرض وارد پروژه کارشناسی نشده‌اند.

فایل اصلی \filepath{main.tex} فقط ترتیب اجزای سند را مشخص می‌کند. مشخصات شخصی در فایل \filepath{config/metadata.tex} قرار دارند. انتخاب قلم در \filepath{config/fonts.tex} و صفحه‌آرایی در \filepath{config/style.tex} نگهداری می‌شود. بنابراین برای استفاده معمول لازم نیست کدهای داخلی قالب تغییر کنند.

\section{شروع سریع}
\subsection{Overleaf}
\begin{enumerate}
  \item فایل ZIP را با گزینه \lr{Upload Project} بارگذاری کنید.
  \item در \lr{Menu}، کامپایلر را روی \lr{XeLaTeX} قرار دهید.
  \item فایل اصلی را \filepath{main.tex} انتخاب کنید.
  \item گزینه \lr{Recompile from scratch} را اجرا کنید.
\end{enumerate}

\subsection{ویندوز}
پس از نصب کامل \lr{TeX Live} یا \lr{MiKTeX}، فایل \filepath{build.bat} را اجرا کنید. این فایل چهار مرحله لازم را خودکار انجام می‌دهد.

\subsection{لینوکس و macOS}
\begin{latin}
\begin{lstlisting}
chmod +x build.sh
./build.sh
\end{lstlisting}
\end{latin}

روش جایگزین برای همه سیستم‌ها:
\begin{latin}
\begin{lstlisting}
latexmk -C
latexmk -xelatex main.tex
\end{lstlisting}
\end{latin}

فقط یک بار اجرای \lr{XeLaTeX} کافی نیست. منابع به \lr{Biber} و فهرست‌ها و ارجاع‌ها به چند اجرای \lr{XeLaTeX} نیاز دارند.

\section{ساختار پروژه}
\begin{longtable}{>{\raggedleft\arraybackslash}p{0.36\textwidth}p{0.54\textwidth}}
\toprule
\textbf{مسیر} & \textbf{کاربرد}\\
\midrule
\filepath{main.tex} & ترتیب صفحات و فصل‌ها\\
\filepath{config/metadata.tex} & مشخصات دانشجو و بخش‌های اختیاری\\
\filepath{config/fonts.tex} & قلم‌های رسمی و جایگزین\\
\filepath{config/style.tex} & اندازه‌ها، شماره‌گذاری و ظاهر\\
\filepath{frontmatter/} & صفحات مقدماتی و انگلیسی\\
\filepath{chapters/} & متن فصل‌ها\\
\filepath{appendices/} & پیوست‌های اختیاری\\
\filepath{assets/} & لوگو و تصاویر\\
\filepath{references.bib} & منابع کتاب‌شناختی\\
\bottomrule
\end{longtable}

\section{تغییر اطلاعات پروژه}
برای ویرایش عنوان، نام دانشجو و استاد فقط \filepath{config/metadata.tex} را تغییر دهید:
\begin{latin}
\begin{lstlisting}[language=TeX]
\newcommand{\ProjectTitleFa}{Persian Project Title}
\newcommand{\ProjectTitleEn}{English Project Title}
\newcommand{\StudentNameFa}{Persian Student Name}
\newcommand{\StudentNumber}{Student Number}
\newcommand{\SupervisorNameFa}{Persian Supervisor Name}
\end{lstlisting}
\end{latin}

فقط متن داخل آکولادها تغییر کند. نام فرمان و آکولادها نباید حذف شوند.

\section{فعال و غیرفعال کردن اجزا}
در انتهای \filepath{config/metadata.tex} کلیدهای زیر وجود دارند:
\begin{latin}
\begin{lstlisting}[language=TeX]
\IncludeDedicationtrue
\IncludeAcknowledgementstrue
\IncludeSymbolstrue
\IncludeListOfFigurestrue
\IncludeListOfTablestrue
\IncludeListOfListingstrue
\IncludeAppendixfalse
\IncludeEnglishPagestrue
\end{lstlisting}
\end{latin}

برای حذف یک بخش، \lr{true} به \lr{false} تبدیل شود. اگر فهرست شکل یا جدول خالی است، کلید مربوط به آن غیرفعال شود.

\section{افزودن فصل}
یک فایل جدید مانند \filepath{chapters/chapter6.tex} بسازید:
\begin{latin}
\begin{lstlisting}[language=TeX]
\chapter{Persian Chapter Title}
\label{chap:new}

\section{Persian Section Title}
Persian chapter text.
\end{lstlisting}
\end{latin}

سپس در \filepath{main.tex} اضافه کنید:
\begin{latin}
\begin{lstlisting}[language=TeX]
\include{chapters/chapter6}
\end{lstlisting}
\end{latin}

\section{متن فارسی و لاتین}
فونت‌های قدیمی فارسی همه حروف لاتین را ندارند. هر اصطلاح لاتین در پاراگراف فارسی باید داخل \commandname{lr} نوشته شود:
\begin{latin}
\begin{lstlisting}[language=TeX]
Persian text \lr{PINN} Persian text \lr{PyTorch}.
\end{lstlisting}
\end{latin}

برای پاراگراف کاملاً انگلیسی از محیط \lr{latin} استفاده شود. رعایت این قاعده از نمایش مربع جلوگیری می‌کند.

\section{منابع و ارجاع}
اطلاعات منبع در \filepath{references.bib} ثبت می‌شود. در متن فارسی از \commandname{pcite} استفاده کنید:
\begin{latin}
\begin{lstlisting}[language=TeX]
\pcite{raissi2019pinns}
\end{lstlisting}
\end{latin}

این فرمان شماره و براکت را با قلم لاتین کامل حروف‌چینی می‌کند. در متن کاملاً انگلیسی، \commandname{cite} معمولی مناسب است.

\section{شکل، جدول و رابطه}
\subsection{شکل}
\begin{latin}
\begin{lstlisting}[language=TeX]
\begin{figure}[H]
  \centering
  \includegraphics[width=.75\textwidth]{figure.pdf}
  \caption{Persian figure caption}
  \label{fig:sample}
\end{figure}
\end{lstlisting}
\end{latin}
عنوان شکل پایین آن قرار می‌گیرد.

\subsection{جدول}
\begin{latin}
\begin{lstlisting}[language=TeX]
\begin{table}[H]
  \centering
  \caption{Persian table caption}
  \label{tab:sample}
  \begin{tabular}{lc}
    \toprule
    Method & Error \\
    \midrule
    PINN & 0.01 \\
    \bottomrule
  \end{tabular}
\end{table}
\end{lstlisting}
\end{latin}
عنوان جدول باید بالای آن باشد.

\subsection{رابطه}
\begin{latin}
\begin{lstlisting}[language=TeX]
\begin{equation}
  E = mc^2
  \label{eq:energy}
\end{equation}
\end{lstlisting}
\end{latin}
ارجاع با \commandname{eqref} انجام می‌شود. رابطه از چپ و شماره آن از راست قرار می‌گیرد.

\section{کد برنامه‌نویسی}
کد و عنوان انگلیسی آن باید چپ‌به‌راست باشند:
\begin{latin}
\begin{verbatim}
\begin{latin}
\begin{lstlisting}[caption={\protect\lr{Python example}},
                  label={lst:python}]
print("Hello")
\end{lstlisting}
\end{latin}
\end{verbatim}
\end{latin}

استفاده از \commandname{lr} در عنوان، از وارونه‌شدن آن در فهرست کدها جلوگیری می‌کند.

\section{قلم‌ها و مربع‌های ناخوانا}
ترتیب قلم‌ها در \filepath{config/fonts.tex} مشخص است. B Mitra برای متن، B Nazanin برای بخش، B Titr برای فصل و Times New Roman برای انگلیسی در اولویت‌اند. در نبود آن‌ها قلم آزاد جایگزین می‌شود.

B Mitra نباید برای اعداد و نمادهای ریاضی انتخاب شود، زیرا بعضی نویسه‌های لازم \lr{unicode-persianmath} را ندارد. قالب برای ریاضیات از Amiri یا قلم کامل دیگر استفاده می‌کند. گلوله فهرست‌ها نیز از فونت ریاضی ساخته می‌شود و به B Mitra وابسته نیست.

\section{قواعد اعمال‌شده از دستورالعمل}
\begin{itemize}
  \item کاغذ A4 و حاشیه ۲٫۵ سانتی‌متر؛
  \item متن ۱۴ پوینت با فاصله خط ۱٫۵؛
  \item تورفتگی ۱٫۲۵ سانتی‌متر؛
  \item عنوان فصل ۱۶ پوینت، پررنگ و وسط‌چین؛
  \item عنوان بخش ۱۴ پوینت، پررنگ و مایل؛
  \item عنوان شکل و جدول ۱۰ پوینت؛
  \item صفحات مقدماتی با حروف فارسی و متن اصلی از صفحه ۱؛
  \item شماره‌گذاری فصل‌محور با خط تیره؛
  \item عنوان جدول در بالا و عنوان شکل در پایین؛
  \item چکیده فارسی و انگلیسی و ۵ تا ۱۰ کلیدواژه.
\end{itemize}

رنگ جلد، صحافی و صفحه سفید پس از جلد، الزامات نسخه فیزیکی هستند و باید هنگام تحویل با گروه آموزشی کنترل شوند.

\section{رفع اشکال}
\subsection{فهرست خالی یا علامت سؤال}
کامپایل چندمرحله‌ای کامل نشده است. \filepath{build.bat}، \filepath{build.sh} یا \lr{latexmk} استفاده شود.

\subsection{مربع به‌جای اصطلاح لاتین}
اصطلاح بیرون از \commandname{lr} نوشته شده است یا یک عنوان انگلیسی با فونت فارسی ساخته شده است.

\subsection{مربع در ارجاع}
\lr{Biber} اجرا نشده یا در متن فارسی از \commandname{cite} به‌جای \commandname{pcite} استفاده شده است.

\subsection{خطای B Mitra و U+066A}
نسخه قدیمی \filepath{config/fonts.tex} در حال استفاده است. پروژه جدید در پوشه‌ای تازه باز شود و فایل‌های کمکی قبلی حذف شوند.

\section{کنترل نهایی}
پیش از تحویل، PDF از ابتدا بازسازی شود و موارد زیر کنترل شوند:
\begin{itemize}
  \item نبود مربع، کلید خام منبع و `؟؟`؛
  \item تکمیل فهرست‌ها و منابع؛
  \item ارجاع همه شکل‌ها و جدول‌ها در متن؛
  \item واقعی‌بودن همه نتایج و نمودارها؛
  \item یکسان‌بودن مفهوم چکیده فارسی و انگلیسی؛
  \item هماهنگی نسخه چاپی با آخرین نظر گروه آموزشی.
\end{itemize}

\end{document}
